% ============================================================================
%  Bruise Segmentation -- Consolidated Results Report
%  Compile:  tools\tectonic\tectonic.exe docs\bruise_full_report.tex
%  Self-contained article class (no aaai26.sty needed).
%  Every number is sourced; see Appendix A for the file each table came from.
% ============================================================================
\documentclass[11pt]{article}
\usepackage[a4paper,margin=1in]{geometry}
\usepackage{booktabs}
\usepackage{array}
\usepackage{amsmath,amssymb}
\usepackage{xcolor}
\usepackage{fancyhdr}
\usepackage{enumitem}
\usepackage{longtable}
\usepackage{caption}

\definecolor{secblue}{RGB}{31,73,125}
\definecolor{boxgray}{RGB}{244,246,249}
\definecolor{wingreen}{RGB}{20,110,60}
\definecolor{losered}{RGB}{150,40,40}
\definecolor{rulegray}{RGB}{200,205,212}

\usepackage[colorlinks=true,linkcolor=secblue,urlcolor=secblue,citecolor=secblue]{hyperref}

\usepackage{sectsty}
\sectionfont{\color{secblue}\large}
\subsectionfont{\color{secblue}\normalsize}
\usepackage[most]{tcolorbox}

\pagestyle{fancy}\fancyhf{}
\renewcommand{\headrulewidth}{0.4pt}
\lhead{\small Bruise Segmentation}\rhead{\small Consolidated Results Report}\cfoot{\thepage}

% right-aligned numeric column types, fixed widths so every table lines up
% (text width at a4 + 1in margins is 15.92cm; every spec below is checked against it)
\newcolumntype{n}{>{\raggedleft\arraybackslash}p{1.35cm}}   % standard numeric
\newcolumntype{M}{>{\raggedleft\arraybackslash}p{1.90cm}}   % numeric w/ wide header
\newcolumntype{C}{>{\raggedleft\arraybackslash}p{3.10cm}}   % confidence interval
\newcolumntype{L}{>{\raggedright\arraybackslash}p{4.6cm}}   % model-name column
\newcolumntype{K}{>{\raggedright\arraybackslash}p{3.4cm}}   % model-name, narrow

\setlength{\parindent}{0pt}\setlength{\parskip}{5pt}
\setlength{\tabcolsep}{5pt}
\renewcommand{\arraystretch}{1.12}

\newtcolorbox{keybox}[1][]{colback=boxgray,colframe=secblue,boxrule=0.6pt,arc=2pt,
  left=8pt,right=8pt,top=6pt,bottom=6pt,#1}
\newcommand{\win}[1]{\textcolor{wingreen}{\textbf{#1}}}
\newcommand{\lose}[1]{\textcolor{losered}{#1}}
\newcommand{\tabnote}[1]{\par\vspace{2pt}{\footnotesize\itshape #1}\par}

\begin{document}

% ===========================================================================
\begin{center}
  {\LARGE\bfseries Bruise Segmentation in White-Light Photographs}\\[5pt]
  {\large Consolidated Results Report}\\[6pt]
  {\small Core models \, $\cdot$ \, best-seed final run \, $\cdot$ \, CNN baselines \, $\cdot$ \,
   capacity scaling \, $\cdot$ \, multi-teacher distillation}\\[3pt]
  {\footnotesize 697/134/185 subject-grouped split \ $\cdot$ \ 2-of-3 consensus target \ $\cdot$ \
   28 test subjects \ $\cdot$ \ ITA skin-tone groups}
\end{center}

\vspace{4pt}
{\color{rulegray}\hrule height 1pt}
\vspace{6pt}

\tableofcontents
\vspace{6pt}
{\color{rulegray}\hrule height 0.5pt}

% ===========================================================================
\section{Executive Summary}

\begin{keybox}
\textbf{In one paragraph.} Eight model configurations spanning $1.6$M--$85$M parameters, CNN and
transformer, were trained and scored under one fixed protocol on a 185-image, 28-subject
consensus test set. Their Dice scores cluster tightly ($0.70$--$0.78$ mean), so \textbf{Dice does
not separate architectures at this sample size} --- the minimum detectable effect is
$\approx0.03$--$0.04$ and every architecture-vs-architecture gap is smaller than that. Two things
\emph{do} separate them. First, the \textbf{complete-miss rate}: SegFormers miss $0.0$--$0.5\%$ of
bruises entirely, YOLO26n misses $2$--$7\%$, and CNN baselines $1.8$--$3.2\%$. Second, a
\textbf{validation oracle} shows the B2 and B5 teachers are genuinely complementary
($+0.026$ Dice of per-image headroom, CI clear of zero) --- so the flat Dice is not a hard ceiling
for \emph{combinations} of models. The distillation study built on that finding produces one
\win{REPORTABLE WIN}: reliability-gated multi-teacher distillation into a 3.7M student matches the
teachers' Dice at $0\%$ complete-miss while \win{significantly improving small-bruise recall}
($+0.053$, CI $[+0.021,+0.076]$), and it does so precisely where naive uniform ensembling of the
same two teachers is \lose{INFERIOR} --- isolating the reliability gate, not "more teachers", as
the effective mechanism.
\end{keybox}

\subsection{Results at a glance}

\begin{center}\small
\begin{tabular}{L n n n n}
\toprule
\textbf{Model} & \textbf{Params} & \textbf{Mean} & \textbf{Median} & \textbf{Miss} \\
 & (M) & Dice & Dice & (\%) \\
\midrule
\multicolumn{5}{l}{\textit{Core models --- 3-seed mean (\S4)}}\\
SegFormer-B2 teacher            & 27.35 & 0.7651 & 0.8113 & 0.00 \\
SegFormer-B0 distilled          & 3.71  & 0.7639 & 0.8114 & 0.18 \\
SegFormer-B0 direct             & 3.71  & 0.7600 & 0.8113 & 0.54 \\
YOLO26n-sem direct              & 1.63  & 0.6911 & 0.8021 & 7.57 \\
YOLO26n-sem distilled           & 1.63  & 0.6676 & 0.7658 & 7.57 \\
\addlinespace
\multicolumn{5}{l}{\textit{CNN baselines --- 3-seed mean (\S6)}}\\
U-Net (ResNet-50)               & $\sim$32 & 0.7526 & 0.8260 & 3.24 \\
DeepLabV3+ (ResNet-50)          & $\sim$40 & 0.7453 & 0.8140 & 1.80 \\
\addlinespace
\multicolumn{5}{l}{\textit{Capacity scaling --- 1 seed (\S7)}}\\
SegFormer-B5                    & $\sim$85 & 0.7829 & 0.8276 & 0.00 \\
\addlinespace
\multicolumn{5}{l}{\textit{Distillation study --- 1 seed, locked eval (\S8)}}\\
\win{p3\_adaptive} (B2$+$B5 $\to$ B0) & 3.71 & \win{0.7748} & 0.8172 & 0.00 \\
\bottomrule
\end{tabular}
\end{center}

\tabnote{Rows are not all mutually comparable: the core models and CNN baselines share one
evaluation, the distillation arms share a separate \emph{locked} re-evaluation (\S8.1). Compare
within a block, and use \S8 for the distillation-arm contrasts.}

\subsection{The four things this report establishes}

\begin{enumerate}[leftmargin=1.6em,itemsep=3pt]
  \item \textbf{Compression is free.} A 3.7M distilled SegFormer-B0 ($0.7639$ mean, 3 seeds) is
  statistically indistinguishable from the 27M B2 teacher ($0.7651$) at $7.4\times$ fewer
  parameters and $2\times$ the throughput.
  \item \textbf{Capacity alone gives diminishing returns.} Going from 3.7M to 85M moves median
  Dice from $0.811$ to $0.828$ --- inside the noise band that the other seven configurations
  also occupy.
  \item \textbf{Complete-miss, not Dice, is the discriminating axis.} It is the only accuracy
  quantity whose model-to-model differences exceed the resolution of this test set
  (YOLO$-$B0-distilled $=+6.5$ points, CI $[+2.4,+11.6]$).
  \item \textbf{Complementary teachers still contain recoverable headroom,} and a per-pixel
  reliability gate --- not plain ensembling --- is what recovers it.
\end{enumerate}

% ===========================================================================
\section{Task, Data, and Evaluation Protocol}

\subsection{Task and target}
Pixel-level segmentation of bruises in white-light (WL) forensic photographs, evaluated for equity
across skin tone. This is a \textbf{multi-labeler study}: three annotators independently labeled
each test image, and the segmentation target is their \textbf{2-of-3 majority-vote consensus mask}.
Training uses one annotator's masks; evaluation is against the three-annotator consensus, making
every test number a \emph{cross-annotator generalization} measurement rather than a conventional
held-out score.

\subsection{Split and protocol}

\begin{center}\small
\begin{tabular}{>{\raggedright\arraybackslash}p{3.4cm} >{\raggedright\arraybackslash}p{11cm}}
\toprule
Split & 697 train / 134 val / 185 test images, subject-grouped, \textbf{zero subject leakage} \\
Test set size & 185 images from \textbf{28 subjects} --- the effective $n$ is 28, not 185 \\
Resolution & $640\times640$ (image bilinear, mask nearest) for all models and all scoring \\
Skin tone & Individual Typology Angle (ITA), 5 groups: Light\,II--III $\to$ Dark\,VI \\
Model selection & best epoch chosen on \textbf{validation} Dice \\
Decision threshold & swept on \textbf{validation} only, then applied \emph{once} to test \\
Seeds & 3 seeds for the core models and CNN baselines; 1 seed for B5 and the distillation arms \\
\bottomrule
\end{tabular}
\end{center}

\subsection{Metrics}

\begin{center}\small
\begin{tabular}{>{\raggedright\arraybackslash}p{4.0cm} >{\raggedright\arraybackslash}p{10.4cm}}
\toprule
Mean / median Dice & Per-image Dice against the consensus mask, then averaged. Median is the more
stable summary; the mean is dragged by the failure tail. \\
\textbf{Complete-miss rate} & Fraction of images where the model predicts \emph{no} bruise on an
image that has one. The high-consequence failure for an injury-documentation tool. \\
Failure-tail metrics & \%\,recall$<$0.10, \%\,Dice$<$0.20, and the 5th-percentile subject-level
Dice. These give resolution where complete-miss saturates at $0\%$. \\
Uncertainty & \textbf{Subject-level cluster bootstrap}, $B=4000$, resampling the 28 subjects.
Model-vs-model comparisons are \textbf{paired} (the same resample scores both models), which is
correct for a shared test set and $\approx2\times$ tighter than independent resampling. \\
$P(\Delta>0)$ & Bootstrap posterior-style probability that the effect is positive. Reported
\emph{alongside} the two-sided CI, never instead of it. \\
\bottomrule
\end{tabular}
\end{center}

\subsection{Inter-annotator context}
On the same 185 test images, pairwise annotator agreement (Dice) ranges $0.58$--$0.75$
(three-pair mean $\approx0.64$), and annotator-vs-consensus ranges $0.70$--$0.87$ (\S5.7). This
indicates substantial human labeling variability, and it frames every model comparison below: the
model-vs-model differences under test ($0.002$--$0.02$) are an order of magnitude smaller than the
human-vs-human disagreement in the labels.

\begin{keybox}
\textbf{Why $n=28$ is a ceiling, not a choice.} There are 28 test subjects because that is how
many subjects \emph{all three} annotators labeled. A 29th cannot be obtained without new
annotation. The correct reading of the null results below is therefore not "this study was too
small" but "\emph{at the maximum sample size for which three-annotator consensus labels exist},
model-level Dice differences are unresolvable."
\end{keybox}

% ===========================================================================
\section{Experiment Inventory}

Every result in this report comes from one of five runs. This table is the map.

\begin{center}\footnotesize
\begin{tabular}{>{\raggedright\arraybackslash}p{2.4cm} >{\raggedright\arraybackslash}p{4.5cm}
                >{\raggedright\arraybackslash}p{1.9cm} >{\raggedright\arraybackslash}p{1.3cm}
                >{\raggedright\arraybackslash}p{3.1cm}}
\toprule
\textbf{Phase} & \textbf{What was run} & \textbf{Where} & \textbf{Seeds} & \textbf{Reported in} \\
\midrule
1 --- Core & SegFormer-B2 / B0-direct / B0-distilled, YOLO26n-sem direct / distilled &
Colab A100 & 3 & \S4 (3-seed), \S5 (best seed) \\
2 --- Baselines & U-Net R50, DeepLabV3+ R50 & Colab A100 & 3 & \S6 \\
3 --- Capacity & SegFormer-B5 & ORC / mlidl & 1 & \S7 \\
4 --- Distillation & 11 KD arms, B2 and/or B5 $\to$ B0 & mlidl 2$\times$A100 & 1 & \S8 \\
--- & nnU-Net v2 (2d) native baseline & ORC MIG slice & 1 (fold 0) & pending \\
\bottomrule
\end{tabular}
\end{center}

\tabnote{The \texttt{*\_a0.xxx} alpha-search trial runs are short-epoch hyperparameter probes and
are excluded from every table in this report.}

% ===========================================================================
\section{Phase 1 --- Core Models (3 Seeds)}

\subsection{Training recipe}
All five core models share one protocol. \textbf{SegFormer} (B2 teacher, B0 student): Dice$+$BCE
loss; AdamW with a backbone/head learning-rate split ($6\times10^{-5}$ / $6\times10^{-4}$), no
weight decay on norm/bias parameters; 1\% warmup then poly decay stepped per optimizer step; AMP
with gradient clipping at 1.0; \textbf{effective batch 8} enforced for every model by gradient
accumulation, so batch size is not confounded with model size. \textbf{YOLO26n-sem}: trained with
the native Ultralytics loop (mosaic, EMA, letterbox, its own schedule and auto-batch), because
forcing a transformer's learning rate onto a CNN detector is not a fair baseline.

\textbf{Distillation} is calibrated response-based KD --- temperature-scaled soft targets:
\[
  \mathcal{L}\;=\;\alpha\,\mathrm{DiceBCE}(z_s,y)\;+\;(1-\alpha)\,\mathrm{BCE}\!\left(z_s,\ \sigma(z_t/T)\right)
\]
with $\alpha=0.6$ for SegFormer-B0 and offline pseudo-mask KD at $\alpha=0.4$ for YOLO
(Ultralytics only consumes a binary mask, so $\alpha>0.5$ would degenerate to the hard label).

\subsection{Three-seed results}

\begin{center}\small
\begin{tabular}{L M M M n}
\toprule
\textbf{Model} & \textbf{Params} & \textbf{Mean Dice} & \textbf{Miss (\%)} & \textbf{Median} \\
 & (M) & (mean $\pm$ s.d.) & (mean $\pm$ s.d.) & Dice \\
\midrule
SegFormer-B2 teacher    & 27.35 & $0.7651\pm0.0037$ & $0.00\pm0.00$ & 0.8113 \\
SegFormer-B0 distilled  & 3.71  & $0.7639\pm0.0047$ & $0.18\pm0.31$ & 0.8114 \\
SegFormer-B0 direct     & 3.71  & $0.7600\pm0.0055$ & $0.54\pm0.00$ & 0.8113 \\
YOLO26n-sem direct      & 1.63  & $0.6911\pm0.0197$ & $7.57\pm1.08$ & 0.8021 \\
YOLO26n-sem distilled   & 1.63  & $0.6676\pm0.0648$ & $7.57\pm4.80$ & 0.7658 \\
\bottomrule
\end{tabular}
\end{center}

\tabnote{YOLO is reported on its \textbf{native-argmax} decoding, which is the sole correct
reporting path (\S5.6). Values are means over seeds 0, 1, 2.}

\subsection{Findings}

\begin{enumerate}[leftmargin=1.6em,itemsep=3pt]
  \item \textbf{Distillation preserves the teacher at $7.4\times$ fewer parameters.} Distilled B0
  ($0.7639$) is indistinguishable from the B2 teacher ($0.7651$) and above direct B0 ($0.7600$).
  Both gaps are inside the $\approx0.03$ minimum detectable effect, so the honest claim is
  \emph{"the student loses nothing measurable"}, not \emph{"distillation helps"}.

  \item \textbf{Distillation acts as a stabiliser more than a booster.} Its clearest effect on
  SegFormer is on \emph{variance and misses}, not the mean: distilled B0's complete-miss rate
  ($0.18\%$) is a third of direct B0's ($0.54\%$), and its seed spread is comparable while its
  worst-case tail is better.

  \item \textbf{Complete-miss separates the architectures where Dice cannot.} All SegFormers miss
  $\le0.5\%$ of bruises entirely; YOLO misses $\approx7.6\%$. YOLO's respectable \emph{median}
  Dice ($0.8021$) actively hides this --- the median is computed over the images it does find.

  \item \textbf{YOLO's misses are seed-unstable.} Across seeds, YOLO-distilled's complete-miss
  rate swings from $2.16\%$ to $11.35\%$ (s.d.\ $4.80$), whereas SegFormer-B2 blanks on
  $0/185$ images in all three seeds. This instability, not the Dice mean, is YOLO's real defect.
\end{enumerate}

\subsection{Per-seed detail}

\begin{center}\small
\begin{tabular}{L n n n n n n}
\toprule
& \multicolumn{3}{c}{\textbf{Mean Dice}} & \multicolumn{3}{c}{\textbf{Complete-miss (\%)}} \\
\cmidrule(lr){2-4}\cmidrule(lr){5-7}
\textbf{Model} & seed 0 & seed 1 & seed 2 & seed 0 & seed 1 & seed 2 \\
\midrule
SegFormer-B2 teacher    & 0.7692 & 0.7638 & 0.7622 & 0.00 & 0.00 & 0.00 \\
SegFormer-B0 direct     & 0.7663 & 0.7579 & 0.7559 & 0.54 & 0.54 & 0.54 \\
SegFormer-B0 distilled  & 0.7680 & 0.7648 & 0.7588 & 0.00 & 0.00 & 0.54 \\
YOLO26n direct          & 0.7028 & 0.6684 & 0.7021 & 8.65 & 7.57 & 6.49 \\
YOLO26n distilled       & 0.7261 & 0.6787 & 0.5978 & 2.16 & 9.19 & 11.35 \\
\bottomrule
\end{tabular}
\end{center}

\tabnote{Native-argmax decoding. The SegFormer seed spread is $\le0.010$ Dice; YOLO-distilled's
spans $0.128$ --- a wider range than \emph{any} architecture difference in this report.}

% ===========================================================================
\section{Phase 1 (continued) --- The Final Run: Best Seed Selected on Validation}
\label{sec:bestseed}

The 3-seed means above answer "how does this architecture behave on average". A deployed system
ships \emph{one} checkpoint. This section reports that single-checkpoint result, chosen the only
legitimate way: \textbf{on validation}.

\subsection{The selection rule}

For each of the five core models:
\begin{enumerate}[leftmargin=1.6em,itemsep=1pt]
  \item Train 3 seeds. For each seed, pick the best epoch on validation Dice.
  \item Sweep the decision cut on \textbf{validation only}; record the val mean Dice at that cut.
  \item \textbf{Select the seed with the highest validation mean Dice.} The test set plays no part
  in this choice.
  \item Run inference \emph{once} with that seed's checkpoint at that seed's val-fitted cut, and
  score on the 185-image test set.
\end{enumerate}

Test inference for this section was re-run from the checkpoints rather than read from a cached
CSV, so the numbers below are recomputed end-to-end and not transcribed.

\subsection{Selection outcome}

\begin{center}\small
\begin{tabular}{L M M M M}
\toprule
\textbf{Model} & \textbf{Selected} & \textbf{Val-fitted} & \textbf{Val mean} & \textbf{Test mean} \\
 & \textbf{seed} & cut & Dice & Dice \\
\midrule
SegFormer-B2 teacher    & 0 & $-0.725$ & 0.7717 & 0.7692 \\
SegFormer-B0 direct     & 0 & $-0.025$ & 0.7558 & 0.7663 \\
SegFormer-B0 distilled  & 0 & $+0.100$ & 0.7662 & 0.7680 \\
YOLO26n direct (native) & 2 & argmax   & 0.7311 & 0.7021 \\
YOLO26n distilled (native) & 0 & argmax & 0.7428 & 0.7261 \\
\bottomrule
\end{tabular}
\end{center}

\tabnote{SegFormer cuts are logits (the sweep is on the raw logit scale, not on probabilities);
native YOLO uses Ultralytics' own argmax and has no swept cut. \textbf{Sanity check that the
selection is honest:} for YOLO-direct, validation picked seed 2, whose test mean ($0.7021$) is
\emph{not} the best of the three seeds (seed 0 reaches $0.7028$). A val-based selector that never
peeked is expected to miss the test optimum occasionally, and here it did.}

\subsection{Full test metrics for the selected checkpoints}

\begin{center}\small
\begin{tabular}{L n n n n n n}
\toprule
\textbf{Model} & \textbf{Mean} & \textbf{Median} & \textbf{Mean} & \textbf{Mean} & \textbf{Mean} &
\textbf{Miss} \\
 & Dice & Dice & IoU & Prec. & Recall & (\%) \\
\midrule
SegFormer-B2 teacher       & 0.7692 & \textbf{0.8192} & 0.6495 & 0.8592 & 0.7281 & \textbf{0.00} \\
SegFormer-B0 distilled     & 0.7680 & 0.8167 & 0.6466 & 0.8625 & 0.7231 & \textbf{0.00} \\
SegFormer-B0 direct        & 0.7663 & 0.8129 & 0.6454 & 0.8599 & 0.7270 & 0.54 \\
YOLO26n distilled (native) & 0.7261 & 0.8012 & 0.6083 & 0.8647 & 0.6777 & 2.16 \\
YOLO26n direct (native)    & 0.7021 & 0.8061 & 0.5939 & \textbf{0.8783} & 0.6697 & \lose{6.49} \\
\bottomrule
\end{tabular}
\end{center}

\tabnote{185 images, 28 subjects, scored at $640\times640$ against the 2-of-3 consensus mask.
Note the shape of YOLO's error: its \emph{precision} is the highest of all five models while its
\emph{recall} is the lowest --- it under-detects rather than over-predicts, which is exactly the
failure mode that produces complete misses.}

\subsection{Failure-tail metrics}

Complete-miss saturates at $0\%$ for the strong SegFormers, so the tail metrics are what give
resolution among them.

\begin{center}\small
\begin{tabular}{L M M M M}
\toprule
\textbf{Model} & \textbf{Miss} & \textbf{\% recall} & \textbf{\% Dice} & \textbf{5th-pct} \\
 & (\%) & $<0.10$ & $<0.20$ & subj.\ Dice \\
\midrule
SegFormer-B2 teacher       & 0.00 & 0.00 & 0.54 & 0.5584 \\
SegFormer-B0 distilled     & 0.00 & 0.00 & 1.08 & \textbf{0.5883} \\
SegFormer-B0 direct        & 0.54 & 1.08 & 1.62 & 0.5634 \\
YOLO26n distilled (native) & 2.16 & 3.78 & 3.78 & 0.4901 \\
YOLO26n direct (native)    & 6.49 & \lose{10.81} & \lose{11.35} & \lose{0.4487} \\
\bottomrule
\end{tabular}
\end{center}

\tabnote{"5th-pct subj.\ Dice" is the 5th percentile of the 28 subject-level mean Dice values ---
i.e.\ how bad the worst subjects are. The two zero-miss SegFormers separate here: B0-distilled has
the better worst-case subject, B2 the better \%\,Dice$<$0.20.}

\subsection{Uncertainty: marginal and paired}

\textbf{Marginal 95\% CIs} (subject-level cluster bootstrap, $B=4000$, seed 42). These overlap
almost completely among the SegFormers --- which is the point.

\begin{center}\small
\begin{tabular}{L M C}
\toprule
\textbf{Model} & \textbf{Mean Dice} & \textbf{95\% CI} \\
\midrule
SegFormer-B2 teacher       & 0.7692 & $[0.7318,\ 0.8024]$ \\
SegFormer-B0 distilled     & 0.7680 & $[0.7387,\ 0.7958]$ \\
SegFormer-B0 direct        & 0.7663 & $[0.7321,\ 0.7982]$ \\
YOLO26n distilled (native) & 0.7261 & $[0.6650,\ 0.7823]$ \\
YOLO26n direct (native)    & 0.7021 & $[0.6271,\ 0.7706]$ \\
\bottomrule
\end{tabular}
\end{center}

\textbf{Paired contrasts.} The same resample scores both models, so these intervals are much
tighter than the marginal ones and are the correct basis for any claim.

\begin{center}\small
\begin{tabular}{>{\raggedright\arraybackslash}p{4.6cm} M C n >{\raggedright\arraybackslash}p{2.85cm}}
\toprule
\textbf{Contrast} & $\Delta$ & \textbf{95\% CI} & $P(\Delta{>}0)$ & \textbf{Verdict} \\
\midrule
Distillation (B0-dist $-$ B0-direct)   & $+0.0017$ & $[-0.0091,+0.0126]$ & 0.582 & n.s. \\
Student $-$ teacher (B0-dist $-$ B2)   & $-0.0012$ & $[-0.0196,+0.0176]$ & 0.437 & n.s. \\
YOLO distilled $-$ direct              & $+0.0239$ & $[-0.0142,+0.0590]$ & 0.901 & n.s. \\
SegFormer $-$ YOLO (B0-dist $-$ YOLO-d)& $+0.0659$ & $[+0.0111,+0.1209]$ & 0.994 & \win{SIGNIFICANT} \\
\textbf{Miss\%}: YOLO-dir $-$ B0-dist & $+6.49$ & $[+2.39,+11.63]$ & 1.000 & \win{SIGNIFICANT} \\
\bottomrule
\end{tabular}
\end{center}

\tabnote{Rows 1--4 are mean Dice; the last row is percentage points of complete-miss rate.
$P(\Delta>0)$ is one-sided and is reported only alongside the two-sided CI.}

\begin{keybox}
\textbf{Reading the nulls correctly.} "n.s." means different things in different rows.
Student$-$teacher at $P=0.437$ is a genuine coin flip: the 3.7M student and the 27M teacher are
interchangeable on this test set. Distillation at $P=0.582$ is likewise unresolved. But
SegFormer$-$YOLO and the miss-rate contrast are \emph{not} null --- and the miss-rate effect
($+6.5$ points, CI excluding zero by a wide margin) is an order of magnitude larger, relative to
its own uncertainty, than anything Dice can see.
\end{keybox}

\subsection{YOLO decoding: native argmax vs the custom \texttt{/255} path}

YOLO26n was scored two ways on the same 185 images and the same selected checkpoints: (a)
\textbf{native argmax}, using Ultralytics' own letterbox preprocessing and class-map argmax; and
(b) a \textbf{custom \texttt{/255}} path that reproduces SegFormer's $640\times640$ stretch
geometry and applies a val-swept probability threshold.

\begin{center}\small
\begin{tabular}{L n n n n n}
\toprule
\textbf{Model / path} & \textbf{Mean} & \textbf{Median} & \textbf{Miss} & \textbf{Mean} &
\textbf{Mean} \\
 & Dice & Dice & (\%) & Prec. & Recall \\
\midrule
YOLO direct --- native argmax     & 0.7021 & 0.8061 & 6.49 & 0.8783 & 0.6697 \\
YOLO direct --- custom \texttt{/255}   & 0.7136 & 0.7934 & 4.32 & 0.8490 & 0.6960 \\
\addlinespace
YOLO distilled --- native argmax  & 0.7261 & 0.8012 & 2.16 & 0.8647 & 0.6777 \\
YOLO distilled --- custom \texttt{/255} & 0.6932 & 0.7674 & 1.08 & 0.7406 & 0.7325 \\
\bottomrule
\end{tabular}
\end{center}

\textbf{Why native argmax is the reported path.} YOLO26n is trained by Ultralytics' own loop,
which feeds letterboxed \texttt{/255} pixels with no ImageNet normalisation. Any evaluation
harness that changes the input distribution reaching the first convolution degrades the logits,
and no downstream threshold or temperature sweep can undo that. Native argmax is the decoding the
model was actually trained for; it is therefore the canonical number, and the custom path is shown
here only to quantify the harness sensitivity.

\begin{keybox}
\textbf{This is itself a result.} The same frozen checkpoint moves by up to $0.033$ mean Dice and
$2.2$ points of miss rate depending only on the evaluation harness --- larger than the
distillation effect ($+0.0017$) and larger than the student--teacher gap ($-0.0012$) that this
project spent most of its compute trying to measure. \textbf{Benchmark preprocessing choices swing
the reported number more than the architecture choices do.}
\end{keybox}

\subsection{The annotation ceiling, and the model that beats its own annotator}
\label{sec:ceiling}

The individual annotator masks allow the labels themselves to be measured on the same 185 images,
with the same code and the same $640$ grid. No model is involved.

\begin{center}\small
\begin{tabular}{>{\raggedright\arraybackslash}p{6.4cm} >{\raggedright\arraybackslash}p{4.6cm} M}
\toprule
\textbf{Comparison} & \textbf{Kind} & \textbf{Mean Dice} \\
\midrule
Gbarimah $\leftrightarrow$ Erik      & human vs human     & 0.7549 \\
Paul $\leftrightarrow$ Erik          & human vs human     & 0.5812 \\
Paul $\leftrightarrow$ Gbarimah      & human vs human     & 0.5809 \\
\quad \textit{average human--human}  &                    & \textit{0.6390} \\
\addlinespace
Gbarimah vs majority                 & human vs consensus & 0.8729 \\
Erik vs majority                     & human vs consensus & 0.8657 \\
\textbf{Paul vs majority}            & human vs consensus & \textbf{0.6998} \\
\addlinespace
SegFormer-B2 teacher                 & model vs consensus & 0.7692 \\
SegFormer-B0 distilled               & model vs consensus & 0.7680 \\
SegFormer-B0 direct                  & model vs consensus & 0.7663 \\
YOLO26n distilled (native)           & model vs consensus & 0.7261 \\
YOLO26n direct (native)              & model vs consensus & 0.7021 \\
\bottomrule
\end{tabular}
\end{center}

Two structural facts follow. First, \textbf{Gbarimah and Erik agree with each other far more than
either agrees with Paul} ($0.755$ vs $0.581$), so the 2-of-3 majority vote is effectively "what
Gbarimah and Erik agreed on". Second, \textbf{Paul --- whose masks every model in this report is
trained on --- matches that consensus target least well of the three annotators} ($0.6998$).

That sets up the comparison: does a model trained only on Paul's masks agree with the three-expert
consensus better than Paul himself does?

\begin{center}\small
\begin{tabular}{>{\raggedright\arraybackslash}p{3.6cm} M M M n >{\raggedright\arraybackslash}p{2.85cm}}
\toprule
\textbf{Model} & \textbf{Model vs} & \textbf{Paul vs} & $\Delta$ & $P(\Delta{>}0)$ &
\textbf{Verdict} \\
 & consensus & consensus & & & \\
\midrule
SegFormer-B2 teacher   & 0.7692 & 0.6998 & $+0.0694$ & 0.990 & \win{SIGNIFICANT} \\
SegFormer-B0 distilled & 0.7680 & 0.6998 & $+0.0682$ & 0.980 & \win{SIGNIFICANT} \\
SegFormer-B0 direct    & 0.7663 & 0.6998 & $+0.0665$ & 0.982 & \win{SIGNIFICANT} \\
YOLO26n distilled      & 0.7261 & 0.6998 & $+0.0263$ & 0.858 & n.s. \\
YOLO26n direct         & 0.7021 & 0.6998 & $+0.0023$ & 0.514 & n.s. \\
\bottomrule
\end{tabular}
\end{center}

\tabnote{Paired subject-level cluster bootstrap, $B=4000$, both sides scored against the same
2-of-3 majority mask. 95\% CIs: B2 $[+0.010,+0.124]$, B0-distilled $[+0.003,+0.133]$,
B0-direct $[+0.005,+0.126]$.}

\begin{keybox}
\textbf{All three SegFormers agree with the three-expert consensus significantly better than the
annotator they were trained on.} Training averages away one annotator's idiosyncrasies and keeps
what generalizes. This effect ($\approx+0.07$) is roughly \textbf{40$\times$ larger than the
distillation effect} ($+0.0017$) --- and unlike it, its confidence interval excludes zero.

It also explains the nulls elsewhere in this report mechanically rather than apologetically:
the project was trying to resolve $0.002$--$0.02$ Dice differences between models against a target
that carries $\approx0.36$ of annotation disagreement ($1-0.639$). \textbf{These models are
label-limited, not capacity-limited.}
\end{keybox}

\subsection{Fairness across ITA skin-tone groups}

Per-group median Dice for the selected checkpoints, with the fairness gap
(best group $-$ worst group) and a Kruskal--Wallis test across the five groups.

\begin{center}\footnotesize
\begin{tabular}{K n n n n n n n}
\toprule
& \multicolumn{5}{c}{\textbf{Median Dice by ITA group}} & & \\
\cmidrule(lr){2-6}
\textbf{Model} & Light & Interm. & Tan & Brown & Dark & \textbf{Gap} & \textbf{KW $p$} \\
 & (39) & (38) & (24) & (29) & (55) & & \\
\midrule
SegFormer-B2 teacher       & 0.789 & 0.867 & 0.773 & 0.810 & 0.793 & 0.094 & \lose{0.028} \\
SegFormer-B0 direct        & 0.827 & 0.834 & 0.810 & 0.813 & 0.770 & 0.064 & 0.709 \\
SegFormer-B0 distilled     & 0.829 & 0.819 & 0.795 & 0.822 & 0.786 & \textbf{0.043} & 0.749 \\
YOLO26n direct (native)    & \lose{0.714} & 0.828 & 0.778 & 0.845 & 0.787 & \lose{0.131} & 0.212 \\
YOLO26n distilled (native) & 0.758 & 0.827 & 0.813 & 0.833 & 0.772 & 0.075 & 0.063 \\
\bottomrule
\end{tabular}
\end{center}

\tabnote{Group sizes (images) in parentheses. Only SegFormer-B2 reaches nominal significance
($p=0.028$), and that is uncorrected across seven model variants.}

The pattern worth naming is in \emph{recall and misses}, not median Dice:

\begin{center}\small
\begin{tabular}{L n n n n}
\toprule
& \multicolumn{2}{c}{\textbf{Light (II--III)}} & \multicolumn{2}{c}{\textbf{All other groups}} \\
\cmidrule(lr){2-3}\cmidrule(lr){4-5}
\textbf{Model} & mean recall & miss (\%) & mean recall & miss (\%) \\
 & ($n=39$) & & ($n=146$) & \\
\midrule
SegFormer-B2 teacher       & 0.734 & 0.0  & 0.726 & 0.0 \\
SegFormer-B0 distilled     & 0.710 & 0.0  & 0.727 & 0.0 \\
SegFormer-B0 direct        & 0.741 & 2.6  & 0.723 & 0.0 \\
YOLO26n direct (native)    & \lose{0.516} & \lose{15.4} & 0.711 & 4.1 \\
YOLO26n distilled (native) & \lose{0.581} & \lose{7.7}  & 0.704 & 0.7 \\
\bottomrule
\end{tabular}
\end{center}

\textbf{Interpretation and the honest caveat.} YOLO under-detects on light skin: its light-skin
recall is $0.52$ against $0.71$ elsewhere, with a $15.4\%$ complete-miss rate in that group, while
its light-skin \emph{precision} ($0.910$) is \emph{above} its own overall mean ($0.878$) ---
under-detection, not confusion. The
SegFormers show no comparable pattern. However, this is \textbf{exploratory}: the Light group has
39 images from only 9 subjects, the fairness gap is a max-minus-min statistic over five groups
(an extremum-of-extrema, which is volatile at this sample size), and there is a real size
confound --- light-skin bruises are smaller in this dataset (median $8{,}085$ px vs $11$--$14$k,
Kruskal $p=0.0003$), and smaller bruises are missed more by every model. The deficit survives
conditioning on size and subject clustering, but only at the boundary
($p\approx0.045$--$0.047$ clustered, versus $p=0.0016$ unclustered --- itself another example of
how much the analysis choice moves the answer). \textbf{No fairness claim in this report should be
stated as established.}

\subsection{Deployment cost}

A100, $640\times640$, batch 1, fp32, 185 images $\times$ 3 repeats, preprocessing excluded from
the timer (raw forward pass only).

\begin{center}\small
\begin{tabular}{L n n n n n}
\toprule
\textbf{Model} & \textbf{Params} & \textbf{Median} & \textbf{p95} & \textbf{FPS} &
\textbf{Peak act.} \\
 & (M) & (ms) & (ms) & & (MB) \\
\midrule
SegFormer-B2 teacher   & 27.35 & 33.67 & 34.95 & 29.7  & 1779 \\
SegFormer-B0 direct    & 3.71  & 16.68 & 16.94 & 59.9  & 1204 \\
SegFormer-B0 distilled & 3.71  & 16.64 & 17.07 & 60.1  & 1204 \\
YOLO26n direct         & 1.63  & 8.18  & 8.36  & 122.2 & 967 \\
YOLO26n distilled      & 1.63  & 8.22  & 8.51  & 121.7 & 967 \\
\bottomrule
\end{tabular}
\end{center}

\begin{keybox}
\textbf{The compression case, stated precisely.} SegFormer-B0-distilled runs $2.02\times$ faster
than the B2 teacher on $7.4\times$ fewer parameters, with $0\%$ complete-miss on both, and a
mean-Dice difference of $-0.0012$ whose 95\% CI is $[-0.020,+0.018]$. There is no measurable
accuracy cost to the compression, and this is the checkpoint to deploy among the five core models.
\end{keybox}

% ===========================================================================
\section{Phase 2 --- Classical CNN Baselines (3 Seeds)}

U-Net and DeepLabV3+ (ResNet-50 encoders, from the
\texttt{segmentation\_models\_pytorch} library) place the transformers against standard
segmentation CNNs under the \emph{identical} split, resolution, training loop, and evaluation.

\begin{center}\small
\begin{tabular}{L M M M n}
\toprule
\textbf{Model} & \textbf{Params} & \textbf{Mean Dice} & \textbf{Miss (\%)} & \textbf{Median} \\
 & (M) & (mean $\pm$ s.d.) & (mean $\pm$ s.d.) & Dice \\
\midrule
U-Net (ResNet-50)      & $\sim$32$^\dagger$ & $0.7526\pm0.0154$ & $3.24\pm0.54$ & \textbf{0.8260} \\
DeepLabV3+ (ResNet-50) & $\sim$40$^\dagger$ & $0.7453\pm0.0182$ & $1.80\pm0.62$ & 0.8140 \\
\bottomrule
\end{tabular}
\end{center}

\tabnote{$^\dagger$ Nominal encoder-based parameter estimates; not measured in this project's
benchmark harness (unlike the SegFormer/YOLO counts in \S5.9, which are).}

\textbf{Finding.} The CNN baselines are competitive on Dice --- U-Net's median ($0.8260$) is
actually the highest of \emph{all} pre-distillation models in this report --- but they miss more
bruises entirely ($1.8$--$3.2\%$) than any SegFormer ($\le0.5\%$). Point gaps to the SegFormers
($0.007$--$0.020$ mean Dice) sit at or below the $\approx0.03$--$0.04$ minimum detectable effect,
so \textbf{this report does not claim SegFormer beats U-Net or DeepLabV3+ on accuracy.} What it
claims is that the three SegFormers and the two CNN baselines now cluster in a
$0.745$--$0.765$ mean-Dice band, and that the
one axis on which they visibly order themselves is complete-miss:
SegFormer $<$ DeepLabV3+ $<$ U-Net $<$ YOLO.

Baseline fairness (ITA) was computed and is not significant for either model.

% ===========================================================================
\section{Phase 3 --- Capacity Scaling: SegFormer-B5}

SegFormer-B5 ($\sim$85M, the largest MiT encoder) was trained with the core SegFormer recipe
verbatim, to test directly whether more capacity raises the ceiling.

\begin{center}\small
\begin{tabular}{L M M n n}
\toprule
\textbf{Model} & \textbf{Params} & \textbf{Mean Dice} & \textbf{Median} & \textbf{Miss (\%)} \\
 & (M) & & Dice & \\
\midrule
SegFormer-B5 (1 seed) & $\sim$85 & 0.7829 & 0.8276 & 0.00 \\
\bottomrule
\end{tabular}
\end{center}

\begin{keybox}
\textbf{Observation across Phases 1--3.} Across eight model configurations spanning $1.6$M--$85$M
parameters, CNN and transformer, median Dice sits in a narrow band ($0.766$--$0.828$): the 85M B5
($0.828$) barely exceeds the 3.7M distilled B0 ($0.811$) and the U-Net baseline ($0.826$).

\emph{Increasing model capacity alone gives diminishing returns under the current data and
evaluation protocol.} This is stated as an observation, not a proven ceiling --- and \S8.2 shows
the qualification matters: a \emph{combination} of models still has measurable headroom that no
single model reaches.
\end{keybox}

\textbf{Status caveat:} B5 is a \textbf{single seed}. Given that the SegFormer seed spread in
\S4.4 is $\pm0.005$--$0.010$, B5's $0.7829$ should be read as "$\approx0.78$, consistent with the
band", not as a rank-1 finish. The 3-seed run is pending, as is the nnU-Net v2 native baseline
(fold 0 trained on an ORC MIG slice; not yet scored).

% ===========================================================================
\section{Phase 4 --- Multi-Teacher Distillation Study}

\subsection{Design and evaluation lock}

Because single-model Dice is flat, the study is framed around a practical question:
\textbf{can a compact 3.7M student preserve overall quality while reducing high-consequence
failures and deployment cost?} It distils into SegFormer-B0 from the best-of-3-seed B5 (selected
on validation, calibrated to $T=2.20$) and/or the B2 teacher.

All KD terms extend the calibrated response loss
\[
  \mathcal{L}\;=\;\alpha\,\mathrm{DiceBCE}(z_s,y)\;+\;(1-\alpha)\,\mathrm{BCE}\!\left(z_s,\ \sigma(z_t/T)\right)
\]
with $\alpha$ searched on validation (Optuna, falling back to grid search). \textbf{Every model in this
phase --- including the B2 and B5 teachers and the B2$\to$B0 reference --- is re-scored under one
locked evaluation}, so all rows in \S8.5 are mutually comparable even though they are not directly
comparable to the Phase 1--3 tables.

\subsection{The multi-teacher gate: a validation oracle}

Before committing compute to multi-teacher methods, a subject-cluster bootstrap on
\textbf{validation} tested whether B2 and B5 are actually complementary. The test set was not
touched.

\begin{center}\small
\begin{tabular}{>{\raggedright\arraybackslash}p{8.6cm} C}
\toprule
\textbf{Quantity (validation only)} & \textbf{Value} \\
\midrule
Spearman(B2 per-image Dice, B5 per-image Dice)          & 0.82 \\
Images where B5 beats B2 by $>0.10$                     & 18 \\
Images where B2 beats B5 by $>0.10$                     & 12 \\
Oracle (best-teacher-per-image) gain over best single    & $+0.0258$ \\
\quad 95\% CI                                            & $[0.0165,\ 0.0362]$ \\
\quad $P(\text{gain}>0)$                                 & 1.00 \\
\textbf{Adaptive gate}                                   & \win{PASS} \\
Reliability weight $\mathrm{rel}_b$ (weight on B5)       & 0.634 \\
\bottomrule
\end{tabular}
\end{center}

\textbf{Interpretation.} B2 and B5 make genuinely different errors on genuinely different images.
A perfect per-image selector would gain $\approx2.6$ Dice points over either teacher alone, with
the CI clear of zero. \textbf{So the flat single-model Dice in \S7 is not a hard ceiling --- it is
a ceiling on \emph{single} models}, and there is measured, real headroom for a method that can
route between complementary teachers.

\subsection{Methods under test}

\begin{center}\small
\begin{tabular}{>{\raggedright\arraybackslash}p{4.6cm} >{\raggedright\arraybackslash}p{9.8cm}}
\toprule
\textbf{Arm} & \textbf{What it transfers} \\
\midrule
Response KD & calibrated teacher probabilities (the baseline KD) \\
CWD (ICCV'21) & per-channel spatial distributions --- the canonical seg-KD baseline \\
BPKD (WACV'24) & boundary and body knowledge, separately \\
Angular (WACV'24) & feature \emph{direction} rather than magnitude (scale-robust for B5$\to$B0) \\
DKD (CVPR'22) & decoupled target / non-target logits (low priority in binary segmentation) \\
Uniform ensemble & $\tfrac12(p_{B2}+p_{B5})$ --- the "more teachers" control \\
\textbf{Reliability-gated (novel)} & per-pixel confidence $\times$ disagreement gate, biased by the
val-derived $\mathrm{rel}_b$ \\
\quad $+$boundary / $+$hard / $+$group / $+$full & one component added at a time (ablation) \\
\bottomrule
\end{tabular}
\end{center}

Teacher mixing itself is not novel (cf.\ HeteroAKD, AAAI'25); the contribution is the specific
combination of a validation-derived reliability gate with failure-risk weighting and boundary
knowledge.

\subsection{Pre-registered success hierarchy}

Because Dice is flat, success is judged in a fixed priority order, decided before the arms ran:

\begin{enumerate}[leftmargin=1.8em,itemsep=1pt]
  \item \textbf{Non-inferiority in subject-level Dice} (margin $0.03$) --- the primary gate.
  \item Reduction in \textbf{complete-miss and the low-recall tail} (\%recall$<$0.10,
        \%Dice$<$0.20, 5th-percentile subject Dice).
  \item Improved \textbf{small-bruise recall}.
  \item Improved \textbf{worst-ITA-group recall} (exploratory).
  \item Lower \textbf{deployment cost} (parameters, latency) --- scored separately.
\end{enumerate}

Every comparison uses the paired subject-cluster bootstrap (28 subjects) against the B2$\to$B0
reference, and each arm receives one verdict:
\win{REPORTABLE WIN} / NON-INFERIOR / \lose{INFERIOR}.

\subsection{Per-arm test results}

All students are SegFormer-B0 (3.71M). Teachers and the reference are shown for context.
Sorted by median Dice.

\begin{center}\footnotesize
\begin{tabular}{>{\raggedright\arraybackslash}p{4.7cm} n n n n n n}
\toprule
\textbf{Model / arm} & \textbf{Mean} & \textbf{Median} & \textbf{Miss} & \textbf{\%rec} &
\textbf{\%Dice} & \textbf{p5} \\
 & Dice & Dice & (\%) & $<$.10 & $<$.20 & subj \\
\midrule
SegFormer-B5 teacher (85M)   & 0.7727 & \textbf{0.8273} & 0.00 & 0.00 & 0.00 & 0.544 \\
\win{p3\_adaptive}           & \win{0.7748} & 0.8172 & 0.00 & 1.08 & 1.08 & 0.623 \\
p2\_cwd                      & 0.7736 & 0.8230 & 0.54 & 2.16 & 2.16 & \textbf{0.650} \\
p3\_adaptive\_boundary       & 0.7668 & 0.8221 & 0.00 & 2.16 & 2.70 & 0.606 \\
SegFormer-B2 teacher (27M)   & 0.7692 & 0.8192 & 0.00 & 0.00 & 0.54 & --- \\
SegFormer-B0 distilled (ref) & 0.7680 & 0.8167 & 0.00 & 0.00 & 1.08 & --- \\
p2\_ensemble\_uniform        & 0.7615 & 0.8139 & 0.00 & 1.62 & 1.62 & 0.515 \\
p2\_bpkd                     & 0.7651 & 0.8091 & 0.54 & 1.08 & 1.62 & 0.589 \\
expA\_response               & 0.7683 & 0.8061 & 0.00 & 2.16 & 2.16 & 0.616 \\
p3\_adaptive\_group          & 0.7586 & 0.8047 & 0.54 & 1.08 & 2.16 & 0.636 \\
x\_angular                   & 0.7418 & 0.8043 & 0.54 & 2.16 & 3.24 & 0.549 \\
p3\_adaptive\_hard           & 0.7414 & 0.8023 & 0.00 & 2.16 & 3.24 & 0.509 \\
p3\_adaptive\_full           & 0.7515 & 0.7984 & 1.08 & 2.70 & 2.70 & 0.587 \\
\bottomrule
\end{tabular}
\end{center}

\tabnote{Locked re-evaluation, 185 images. Single seed (42) per arm. Because the teachers are
re-scored here too, their numbers differ slightly from \S5.3 --- compare within this table only.}

\subsection{Statistical verdicts vs the B2$\to$B0 reference}

\begin{center}\footnotesize
\begin{tabular}{>{\raggedright\arraybackslash}p{3.4cm} >{\raggedright\arraybackslash}p{2.7cm}
                >{\raggedleft\arraybackslash}p{3.4cm} >{\raggedright\arraybackslash}p{4.1cm}}
\toprule
\textbf{Arm} & \textbf{Verdict} & $\Delta$ \textbf{small recall [95\% CI]} & \textbf{Note} \\
\midrule
\win{p3\_adaptive}     & \win{REPORTABLE WIN} & $+0.053\,[+0.021,+0.076]$ &
  non-inferior Dice $+$ small-recall gain \\
p2\_cwd                & NON-INFERIOR & $+0.004\,[-0.053,+0.049]$ & compresses without loss \\
p2\_bpkd               & NON-INFERIOR & $-0.020\,[-0.080,+0.044]$ & compresses without loss \\
p3\_adaptive\_boundary & NON-INFERIOR & $-0.010\,[-0.083,+0.039]$ & \\
expA\_response         & NON-INFERIOR & $-0.028\,[-0.055,-0.007]$ & regresses small recall \\
p2\_ensemble\_uniform  & \lose{INFERIOR} & $-0.002\,[-0.049,+0.026]$ & subject-Dice CI lo $-0.031$ \\
p3\_adaptive\_group    & \lose{INFERIOR} & $-0.006\,[-0.060,+0.021]$ & subject-Dice CI lo $-0.037$ \\
p3\_adaptive\_full     & \lose{INFERIOR} & $-0.076\,[-0.128,-0.018]$ & subject-Dice CI lo $-0.042$ \\
p3\_adaptive\_hard     & \lose{INFERIOR} & $-0.024\,[-0.058,+0.019]$ & subject-Dice CI lo $-0.052$ \\
x\_angular             & \lose{INFERIOR} & $-0.054\,[-0.111,-0.011]$ & subject-Dice CI lo $-0.065$ \\
\bottomrule
\end{tabular}
\end{center}

\tabnote{"$\Delta$ small recall" is (arm $-$ reference) recall on small bruises. INFERIOR verdicts
are driven by the primary gate: the lower bound of the paired subject-Dice CI falls below the
$-0.03$ non-inferiority margin.}

\subsection{Findings}

\begin{enumerate}[leftmargin=1.6em,itemsep=4pt]
  \item \textbf{The reliability-gated adaptive method is the sole winner.} \win{p3\_adaptive} is
  the only \textbf{REPORTABLE WIN}: non-inferior on subject Dice \emph{and} a significant
  improvement in small-bruise recall ($+0.053$, CI $[+0.021,+0.076]$, $P=0.996$) --- a real
  reduction in a high-consequence failure that Dice hides. At $3.71$M parameters its mean Dice
  ($0.7748$) edges out both the 85M B5 and the 27M B2 teachers under the same locked eval, at
  $0\%$ complete-miss.

  \item \textbf{The gate matters, not just "more teachers".} With the \emph{same two teachers},
  plain uniform averaging is \lose{INFERIOR} (subject-Dice CI lower bound $-0.031$), while the
  per-pixel confidence$\times$disagreement gate wins. This contrast is the cleanest evidence in
  the study: it isolates the \emph{gating mechanism}, not the ensemble, as what captures the
  complementary knowledge the validation oracle promised.

  \item \textbf{The plain gate is best; the failure-targeting add-ons over-engineer it.} Every
  ablation stacked on the adaptive base reduced subject Dice --- \lose{$+$hard}, \lose{$+$group},
  and \lose{$+$full} all to INFERIOR, and $+$boundary down to merely non-inferior. The reliability
  gate alone is the effective ingredient.

  \item \textbf{Standard seg-KD baselines compress without loss.} CWD, BPKD, and response KD are
  all non-inferior --- consistent with \S4.3's finding that the compact student loses nothing.
  Angular distillation \lose{hurt} ($0.7418$): matching feature \emph{direction} under-performed
  here, plausibly because the B5$\to$B0 architectural gap makes the feature spaces
  non-corresponding.
\end{enumerate}

\begin{keybox}
\textbf{Phase 4 in one sentence.} A 3.7M student distilled from a reliability-gated fusion of two
complementary teachers matches the teachers' overall Dice at $0\%$ complete-miss while
significantly improving small-bruise recall --- and does so precisely where naive ensembling of
the same teachers fails, identifying the reliability gate as the effective mechanism.
\end{keybox}

% ===========================================================================
\section{Synthesis}

Reading Phases 1--4 together, three claims are supported and one common claim is not.

\textbf{Supported.}
\begin{enumerate}[leftmargin=1.6em,itemsep=2pt]
  \item \textbf{Compression is free on this task.} A 3.7M student matches a 27M teacher
  (paired $\Delta=-0.0012$, CI $[-0.020,+0.018]$) at $2\times$ the throughput, and every standard
  seg-KD method reproduces this. Deploy the small model.
  \item \textbf{Failure rate, not Dice, is the axis that resolves.} Complete-miss ranges $0\%$ to
  $7.6\%$ across the model set, and the SegFormer--YOLO miss contrast ($+6.5$ points, CI
  $[+2.4,+11.6]$) is one of only three significant accuracy findings in the whole project.
  \item \textbf{The labels are the binding constraint.} Human--human agreement is $0.639$; all
  three SegFormers agree with the consensus significantly \emph{better} than the annotator they
  learned from ($\approx+0.07$, $P\ge0.98$). Model-vs-model differences ($0.002$--$0.02$) are far
  below the annotation noise floor.
\end{enumerate}

\textbf{Not supported.} That any architecture beats any other on Dice. The six
non-YOLO configurations span just $0.745$--$0.783$ mean Dice; the minimum detectable effect at $n=28$ subjects is
$\approx0.03$--$0.04$; and, as \S5.6 shows, the evaluation harness alone can move a frozen
checkpoint by $0.033$. Architecture-ranking claims on this test set are not identifiable.

\textbf{The one place a real gain was found} is not a bigger model but a \emph{better use of two
models}: the validation oracle measured $+0.026$ of per-image headroom between B2 and B5, and
reliability-gated distillation converted part of that into a significant small-bruise recall
improvement in a 3.7M student.

% ===========================================================================
\section{Limitations}

\begin{enumerate}[leftmargin=1.6em,itemsep=3pt]
  \item \textbf{Single seed for the distillation arms and for B5.} The subject-cluster bootstrap
  captures test-subject sampling variability, \emph{not} training-seed variance --- so all Phase 3
  and Phase 4 intervals are optimistic. Given that YOLO-distilled's seed spread is $0.128$ Dice
  (\S4.4), this matters. \texttt{p3\_adaptive} and \texttt{p2\_ensemble\_uniform} should be re-run
  at 3 seeds before the Phase 4 conclusion is treated as locked.
  \item \textbf{Margins are modest.} Most Phase 4 differences sit inside the Dice MDE
  ($\approx0.03$--$0.04$); the win is carried by the failure-tail endpoint, exactly as the
  pre-registered hierarchy intended, but it is a tail effect and should be described as one.
  \item \textbf{$n=28$ subjects is a hard ceiling} set by three-annotator overlap, not by budget.
  No amount of additional compute changes the resolution of this test set.
  \item \textbf{Fairness is exploratory.} 9--17 subjects per ITA group; the fairness gap is an
  extremum-of-extrema statistic; a real bruise-size confound is present; and clustered vs
  unclustered standard errors move the light-skin result from $p=0.0016$ to $p\approx0.046$.
  Nothing here is an established fairness finding.
  \item \textbf{Cross-phase comparability.} Phase 1--3 and Phase 4 use different evaluation
  passes (the Phase 4 lock re-scores the teachers). Rows must be compared within a phase.
  \item \textbf{Training pool is one annotator's masks.} Whether the "model beats its annotator"
  result generalises to a model trained on Gbarimah or Erik is untested --- their training pools
  are not on disk.
  \item \textbf{Single site, single camera setup.} No external-dataset validation has been run.
\end{enumerate}

% ===========================================================================
\section{Next Steps}

\begin{enumerate}[leftmargin=1.6em,itemsep=3pt]
  \item \textbf{Three-seed the Phase 4 winner and its control} (\texttt{p3\_adaptive} and
  \texttt{p2\_ensemble\_\allowbreak uniform}). This is the single highest-value remaining run: it converts the
  study's headline from suggestive to locked.
  \item \textbf{Finish the pending baselines}: SegFormer-B5 at 3 seeds, and nnU-Net v2 fold 0
  scored (it is trained but not yet scored). Both slot into the annotation-ceiling table.
  \item \textbf{Distil the winning teacher configuration into YOLO.} YOLO's $\approx7.6\%$
  complete-miss is the largest real failure in the model set and the only one big enough to move
  significantly.
  \item \textbf{Retrain on Gbarimah's and Erik's masks} (equalised pool sizes, same architecture
  and hyperparameters, scored against the same consensus). This directly tests "who labels your
  data matters more than which model you train" --- given \S5.7, it is likely the strongest
  available result and needs no new annotation.
  \item \textbf{Run the harness-sensitivity sweep} (normalisation $\times$ letterbox $\times$
  interpolation on frozen weights). \S5.6 already shows one instance of it; systematising it turns
  a pipeline defect into a general claim about benchmark fragility.
  \item \textbf{Data scaling and re-annotation.} Distillation cannot answer whether performance is
  genuinely data-limited; only adding data or adding annotators can.
\end{enumerate}

% ===========================================================================
\appendix
\section{Provenance of Every Number}

Only results computed in this project are treated as ground truth. Older reference PDFs in
\texttt{docs/} are used for design rationale only, never for numbers.

\begin{center}\footnotesize
\begin{tabular}{>{\raggedright\arraybackslash}p{4.1cm} >{\raggedright\arraybackslash}p{10.3cm}}
\toprule
\textbf{Section} & \textbf{Source} \\
\midrule
\S4.2, \S4.4 (3-seed, per-seed) &
\texttt{results\_final/FINAL\_RESULTS.csv}, \texttt{segformer\_test\_per\_seed.csv},
\texttt{yolo\_test\_per\_seed.csv} \\
\S5.2 (selection outcome) &
\texttt{best\_seed\_val\_selected/best\_seed\_val\_selected\_results.csv} \\
\S5.3--5.4 (test + tail metrics) &
per-image CSVs in \texttt{final\_analysis\_.../per\_image\_*.csv}, 185 rows each; tail metrics
recomputed here from those rows \\
\S5.5 (marginal CIs) &
recomputed for this report: subject-level cluster bootstrap, $B{=}4000$, seed 42, over the same
per-image CSVs \\
\S5.5 (paired contrasts) & \texttt{final\_analysis\_.../paired\_contrasts.csv} \\
\S5.6 (YOLO paths) & \texttt{per\_image\_yolo\_sem\_*\_custom255.csv} vs the native per-image CSVs \\
\S5.7 (annotation ceiling) &
\texttt{interlabeler\_agreement\_640.csv} $\to$ \texttt{annotation\_ceiling.csv},
\texttt{model\_vs\_paul.csv} \\
\S5.8 (fairness) &
\texttt{fairness\_per\_group\_bestseed.csv}, \texttt{fairness\_stats\_bestseed.csv} \\
\S5.9 (speed) & \texttt{results\_final/benchmark\_640.csv} (A100, fp32, batch 1) \\
\S6 (CNN baselines) &
\texttt{results\_baselines/BASELINES\_RESULTS.csv},
\texttt{smp\_baselines\_test\_per\_seed.csv} \\
\S7 (B5) & SegFormer-B5 baseline run, locked-eval score \\
\S8 (distillation) &
distillation-suite aggregate report (mlidl, 2$\times$A100); all arms single seed under one locked
re-evaluation \\
\bottomrule
\end{tabular}
\end{center}

\vspace{6pt}
{\color{rulegray}\hrule height 0.5pt}
\vspace{4pt}

\begin{center}\footnotesize\itshape
Phases 1--2 are 3-seed means (native-argmax for YOLO); SegFormer-B5 and all distillation arms are
single seed. Everything is scored on the 185-image, 28-subject consensus test set at
$640\times640$ under one fixed protocol, with thresholds fitted on validation and applied once to
test. The distillation arms use a separate locked re-evaluation for cross-arm comparability.
The \texttt{*\_a0.xxx} alpha-search trial runs are short-epoch hyperparameter probes and are
excluded from all tables.
\end{center}

\end{document}
